\chapter{第十一届大唐杯本科A组省赛}
\fancyhead[L]{\color{H6}\kaishu\faIcon{atom}\;第十一届大唐杯本科A组省赛}
\fancyhead[R]{\color{H6}\kaishu\rightmark\,}

\date{2024年4月20日}{}{\href{https://github.com/xubenshan/dtcup}{\textbf{才疏学浅的小熊}}}
{}
{https://www.xiaohongshu.com/user/profile/6192219d0000000021028647}{小红书}
{https://space.bilibili.com/692546609}{B站账号}



\section{单选题（共40道题，共120分）}



\begin{choice}{D}[]
    5G切换过程分为测量、判决、执行三个阶段，其中判决结果由哪个网元给出
    \begin{tasks}(4)
        \task 核心网AMF
        \task 目标基站
        \task 终端
        \task 源基站
    \end{tasks}
\end{choice}


\begin{choice}{D}[]
    5G设备EMB6116加装理线架总计高度是
    \begin{tasks}(4)
        \task 3U
        \task 2U
        \task 5U
        \task 4U
    \end{tasks}
\end{choice}




\begin{choice}{A}[]
    常用的固定卫星以C频段的居多，C频段为上行6GHZ、下行多少GHz
    \begin{tasks}(4)
        \task 4
        \task 8
        \task 10
        \task 6
    \end{tasks}
\end{choice}




\begin{choice}{A}[]
    作业成本法认为，将成本分配到成本对象有多少种不同的形式。
    \begin{tasks}(4)
        \task 3
        \task 4
        \task 5
        \task 2
    \end{tasks}
\end{choice}


\begin{choice}{C}[]
    IPMT是指
    \begin{tasks}(4)
        \task 集成决策管理团队
        \task 集成技术管理团队
        \task 集成组合管理团队
        \task 集成管理层团队
    \end{tasks}
\end{choice}

\begin{choice}{A}[]
    C-V2X的工作模式不包括
    \begin{tasks}(4)
        \task V2C
        \task V2V
        \task V2I
        \task V2P
    \end{tasks}
\end{choice}




\begin{choice}{D}[]
    NR系统中，随机接入的目的错误的是
    \begin{tasks}(2)
        \task 获取C-RNTI，基站识别UE的标识
        \task 实现UE和gNB之间的上行同步
        \task 获取MSG3的资源
        \task 实现UE和gNB之间的下行同步
    \end{tasks}
\end{choice}


\begin{choice}{D}[]
    5G系统中终端可通过系统消息获取小区选择相关参数，该系统消息为
    \begin{tasks}(4)
        \task SIB2
        \task SIB3
        \task MIB
        \task SIB1
    \end{tasks}
\end{choice}


\begin{choice}{C}[]
    5G RS232接口标准是一种串行通信的标准，是实现以下哪个的通信方式
    \begin{tasks}(4)
        \task 多对点
        \task 点对多
        \task 点对点
        \task 多对多
    \end{tasks}
\end{choice}


\begin{choice}{B}[]
    5G帧结构中采用常规CP时，一个时隙中固定有多少个符号
    \begin{tasks}(4)
        \task 20
        \task 14
        \task 6
        \task 12
    \end{tasks}
\end{choice}

\begin{choice}{A}[]
    在5GR15版本中，SIB1消息中TDD-UL-DL-ConfigCommon中dI-UL-
    Transmission Periodicity为ms5（6），以及四元组配置为（7，6，2，4），由此可知帧结构5ms中上行时隙数是
    \begin{tasks}(4)
        \task 2
        \task 7
        \task 6
        \task 4
    \end{tasks}
\end{choice}


\begin{choice}{A}[]
    5G网络中，上行流量校验及上行流量上报是哪个网元的功能
    \begin{tasks}(2)
        \task UPF
        \task AMF
        \task AUSF
        \task SMF
    \end{tasks}
\end{choice}


\begin{choice}{C}[]
    5G智能网络优化调参后，最终模型评估时使用的数据集是
    \begin{tasks}(2)
        \task 训练集
        \task 验证集
        \task 测试集
        \task 以上都不是
    \end{tasks}
\end{choice}


\begin{choice}{D}[]
    移动通信中常规天线的极化方式为（\quad）
    \begin{tasks}(2)
        \task 圆极化
        \task 水平极化
        \task 双极化
        \task 垂直极化
    \end{tasks}
\end{choice}

\begin{choice}{B}[]
    卫星互联网通信中，（\quad）的性能更适配卫星互联网发展
    \begin{tasks}(4)
        \task 都可以
        \task 低轨卫星
        \task 中轨卫星
        \task 高轨卫星
    \end{tasks}
\end{choice}

\begin{choice}{B}[]
    卫星移动通信业务使用哪个频段
    \begin{tasks}(4)
        \task C
        \task L
        \task S
        \task Ku
    \end{tasks}
\end{choice}



\begin{choice}{A}[]
    5G NR中PCI数量可以规划（\quad）个
    \begin{tasks}(4)
        \task 1008
        \task 504
        \task 1512
        \task 2016
    \end{tasks}
\end{choice}

\begin{choice}{C}[]
    提供同频小区重选的相关信息的系统消息是
    \begin{tasks}(4)
        \task SIB2
        \task MIB
        \task SIB3
        \task RMSI
    \end{tasks}
\end{choice}


\begin{choice}{B}[]
    5G场景中，智能抄表业务对哪项指标要求较高
    \begin{tasks}(4)
        \task 峰值速率
        \task 连接数
        \task 时延
        \task 速率
    \end{tasks}
\end{choice}



\begin{choice}{B}[]
    我国的固定信令系统No.7信令网采用（\quad）级分类
    \begin{tasks}(4)
        \task 4
        \task 3
        \task 2
        \task 5
    \end{tasks}
\end{choice}



\begin{choice}{B}[]
    未来6G网络将使用分布式大规模MIMO技术，将采用（\quad）的网络架构。
    \begin{tasks}(2)
        \task 高基站密度
        \task 以用户为中心
        \task 集中式大规模
        \task 以基站为中心
    \end{tasks}
\end{choice}




\begin{choice}{D}[]
    如果决策树过度拟合训练集，以下不属于消除过拟合的选项是
    \begin{tasks}(2)
        \task 剪枝
        \task 采用随机森林，集成学习
        \task 减少maxdepth
        \task 增加maxdepth
    \end{tasks}
\end{choice}

\begin{choice}{C}[]
    DC供电系统的电压波动范围为（\quad）
    \begin{tasks}(2)
        \task -40V至-55V
        \task -40V至-50V
        \task -40V至-57V
        \task -40V至-53V
    \end{tasks}
\end{choice}

\begin{choice}{C}[]
    在5GNSA场景中，SCG的主小区被称作
    \begin{tasks}(2)
        \task MCG Secondary cell
        \task SCG Secondary cell
        \task PSCell
        \task Primary cell
    \end{tasks}
\end{choice}


\begin{choice}{D}[]
    关于扩频码和扰码描述不正确的是
    \begin{tasks}(1)
        \task 扩频码需要改变带宽，扰码需要改变信号的带宽
        \task 扩频码不需要改变带宽，扰码不改变信号的带宽
        \task 扩频码不需要改变带宽，扰码不改变信号的带宽
        \task 扩频码需要改变带宽，扰码不改变信号的带宽
    \end{tasks}
\end{choice}


\begin{choice}{A}[]
    下列选项中，哪个是NFV的技术基础
    \begin{tasks}(4)
        \task 虚拟化
        \task 云计算
        \task 人工智能
        \task 大数据
    \end{tasks}
\end{choice}



\begin{choice}{B}[]
    5G中小区异频同优先级或低优先级重选启测条件为
    \begin{tasks}(1)
        \task 始终需要测量
        \task 满足Srxlev<=SnonintraSearchPand \;Squal<=SnonintraSearchQ时需测量
        \task 满足Srxlev<=SintraSearchPand \;Squal<=SintraSearchQ时需测量
        \task 满足Srxlev>0且Squal>0时，需测量
    \end{tasks}
\end{choice}


\begin{choice}{A}[]
    在5G NR中，一个REG（Resource Element Group）包含多少个RB（resource block）？
    \begin{tasks}(4)
        \task 1
        \task 2
        \task 3
        \task 4
    \end{tasks}
\end{choice}


\begin{choice}{D}[]
    IPD模式流程中，哪一阶段才能开始产品的正式立项？
    \begin{tasks}(2)
        \task IPMT第一阶段
        \task IPMT第四阶段
        \task IPMT第三阶段
        \task IPMT第二阶段
    \end{tasks}
\end{choice}


\begin{choice}{B}[]
    在5G NR中，PDCCH采用的调制方式是

    \begin{tasks}(4)
        \task 256QAM
        \task QPSK
        \task 64QAM
        \task 16QAM
    \end{tasks}
\end{choice}


\begin{choice}{B}[]
    华为Mate60Pro的卫星通话功能采用中国自主研制的天通一号卫星移动通信系统目前由几颗通信卫星支撑。
    \begin{tasks}(4)
        \task 2
        \task 3
        \task 4
        \task 5
    \end{tasks}
\end{choice}


\begin{choice}{D}[]
    我国的固定信令系统No.7信令网采用几级分类。
    \begin{tasks}(4)
        \task 5
        \task 2
        \task 1
        \task 3
    \end{tasks}
\end{choice}



\begin{choice}{D}[]
    30W功率折算到dB域为多少dBm
    \begin{tasks}(4)
        \task 30
        \task 36
        \task 46
        \task 45
    \end{tasks}
\end{choice}


\begin{choice}{D}[]
    NR系统中，CSI可以包括CQI、PMI、CRI、SSBRI、LI、RIL以及L1-RSRP。其中，LI用于指示
    \begin{tasks}(2)
        \task 资源指示
        \task 波束索引
        \task 波束强度
        \task PMI中最强的列
    \end{tasks}
\end{choice}


\begin{choice}{C}[]
    PDT计划阶段的主要工作内容不包括
    \begin{tasks}(2)
        \task 新产品市场需求确认
        \task 新产品中试需求确认
        \task 新产品投资回报率确认
        \task 新产品开发模式确认
    \end{tasks}
\end{choice}


\begin{choice}{C}[]
    使用托管交换机替换非托管交换机的一个原因是
    \begin{tasks}(2)
        \task 可减少冲突域
        \task 可在网络之间路由
        \task 可支持多个VLAN
        \task 可管理路由器表
    \end{tasks}
\end{choice}


\begin{choice}{B}[]
    5G帧结构描述中，下面哪一项是错误的
    \begin{tasks}(1)
        \task 帧结构配置可以由SIB静态帧结构配置
        \task R15的协议中，RRC高层配置的tdd-UL-DL-ConfigurationDedicated定义了周期大小
        \task 上下行资源比例可在1：4到2:3之间调整
        \task DCI format2-0用于动态指示帧结构
    \end{tasks}
\end{choice}


\begin{choice}{C}[]
    5G系统中，EPS fallback方案，提供语音业务的网络是
    \begin{tasks}(4)
        \task 5G
        \task GSM
        \task LTE
        \task WCDMA
    \end{tasks}
\end{choice}


\begin{choice}{B}[]
    手机卫星电话通信允许用户在全球范围内进行语音通话，根据卫星电话系统的工作原理，该系统不包括
    \begin{tasks}(2)
        \task 通信卫星组成的链路
        \task 卫星天线设备
        \task 地面站设备
        \task 用户终端设备
    \end{tasks}
\end{choice}


\begin{choice}{A}[]
    产品变动成本管理要做的第一件基础性工作是
    \begin{tasks}(2)
        \task 根据产品型号建立产品变动成本数据库
        \task 建立产品变动成本考核体系
        \task 确定数据库的成本数据项字段
        \task 数据库中产品作业环节变动成本数据录入
    \end{tasks}
\end{choice}



\section{多选题（共30小题，共120分）}
\begin{choice}{\;BCD\;}[]
    随着工业互联网安全攻击日益呈现出的新型化、多样化、复杂化，现有的工业互联网安全暴露出哪些问题
    \begin{tasks}(1)
        \task 工业生产迭代周期短，存量设备可以快速进行安全防护升级换代
        \task 工业信息安全存在先天不足，安全防护能力难以快速提升
        \task OT与IT两个领域人员融合较慢，安全意识急需提升
        \task 数据隐私和数据安全防护缺乏有效手段
    \end{tasks}
\end{choice}

\begin{choice}{\;ABCD\;}[]
    电信网的传输设备根据传输介质的不同分为：
    \begin{tasks}(2)
        \task 无线传输设备
        \task 光纤传输设备、
        \task 微波传输设备
        \task 缆线传输设备
    \end{tasks}
\end{choice}

\begin{choice}{\;BCD\;}[]
    关于激光雷达说法正确的是
    \begin{tasks}(1)
        \task 不受大气和气象限制
        \task 可以获得目标反射的幅度、频率和相位等信息
        \task 全天候工作，不受白天和黑夜光照条件的限制
        \task 抗干扰性能好
    \end{tasks}
\end{choice}

\begin{choice}{\;BCD\;}[]
    如下选项中哪些属于6G网络架构中所定义的“五面”组成部分
    \begin{tasks}(4)
        \task 传输面
        \task 数据面
        \task 控制面
        \task 智能面
    \end{tasks}
\end{choice}


\begin{choice}{\;ABD\;}[]
    决策树学习通常包括哪几个步骤：
    \begin{tasks}(2)
        \task 决策树的生成
        \task   特征选择
        \task 参数估计
        \task 决策树的剪枝
    \end{tasks}
\end{choice}

\begin{choice}{\;AC\;}[]
    NR中T300定时器下面说法正确的是
    \begin{tasks}(1)
        \task 在接收到RRCSetup或RRCReject消息后停止定时器
        \task 在小区重选，并在高层中断连接建立时停止定时器
        \task 在发送RRCSetupRequest时启动定时器
        \task 在接收到RRCReestablishment消息后停止定时器
    \end{tasks}
\end{choice}

\begin{choice}{\;ABC\;}[]
    NSA场景下，大唐4G和5G基站间需要建立ENDC链路。下列关于ENDC链路的建立说法正确的是
    \begin{tasks}(1)
        \task 建立ENDC链路时，需要先建立服务端，再建立客户端
        \task 建立ENDC链路时两端协议类型必须都是ENDC
        \task 如果小区未建立，不会触发ENDC建立流程
        \task ENDC关联的两个基站不需要建立路由关系
    \end{tasks}
\end{choice}


\begin{choice}{\;AB\;}[]
    集成产品开发模式把项目管理分为了哪些方面
    \begin{tasks}(4)
        \task PDT
        \task IPMT
        \task PMP
        \task PMO
    \end{tasks}
\end{choice}


\begin{choice}{\;ABC\;}[]
    5G NR系统PRACH格式支持
    \begin{tasks}(2)
        \task format1
        \task formatA1
        \task formatB1
        \task formatC1
    \end{tasks}
\end{choice}


\begin{choice}{\;BD\;}[]
    5G帧结构中的循环前缀CP类型有
    \begin{tasks}(4)
        \task 特殊CP
        \task 常规CP
        \task 密集CP
        \task 扩展CP
    \end{tasks}
\end{choice}


\begin{choice}{\;ABCD\;}[]
    数字卫星通信方式有

    \begin{tasks}(2)
        \task 星上交换时分多址
        \task 不加话音插空的时分多址
        \task 卫星数字通道方式
        \task 数字话音插空的时分多址
    \end{tasks}
\end{choice}

\begin{choice}{\;AD\;}[]
    在5G网络中，当UE处于RRC-inactive状态，以下描述正确的是
    \begin{tasks}(2)
        \task UE会保留呼叫的上下文
        \task UE进入新的小区会发起切换
        \task 核心网认为UE处于空闲状态
        \task RAN会保留呼叫的上下文
    \end{tasks}
\end{choice}

\begin{choice}{\;ABCD\;}[]
    PDT团队的作用包括哪些方面
    \begin{tasks}(4)
        \task 中试
        \task 销售
        \task 开发
        \task 服务
    \end{tasks}
\end{choice}

\begin{choice}{\;BCD\;}[]
    在5G网络中宏站系统内的邻区添加规则为
    \begin{tasks}(1)
        \task 添加第二圈所有小区为邻区
        \task 添加第一圈所有小区为邻区
        \task 添加第二圈正对的小区为邻区
        \task 添加本站其他小区为邻区
    \end{tasks}
\end{choice}

\begin{choice}{\;ACD\;}[]
    作业成本法认为，将成本分配到成本对象有以下哪几种形式？

    \begin{tasks}(2)
        \task 成本追溯
        \task 资源分解
        \task  分摊
        \task 动因归集
    \end{tasks}
\end{choice}

\begin{choice}{\;BCD\;}[]
    以下哪个环节的成本性态为变动成本？
    \begin{tasks}(2)
        \task 生产管理
        \task 采购管理
        \task 整机测试
        \task 单板测试
    \end{tasks}
\end{choice}

\begin{choice}{\;BCD\;}[]
    手机卫星电话通信允许用户再全球范围内进行语音通话，根据卫星电话系统的工作原理，该系统需要包括如下设备
    \begin{tasks}(2)
        \task 卫星天线设备
        \task 地面站设备
        \task 用户终端设备
        \task 通信卫星组成的链路
    \end{tasks}
\end{choice}

\begin{choice}{\;AD\;}[]
    下面有关5G PRACH功率控制说法正确的是
    \begin{tasks}(2)
        \task 只有开环没有闭环
        \task 全部路径补偿
        \task 开环+闭环功率控制
        \task 每多发送一次会增加功率
    \end{tasks}
\end{choice}

\begin{choice}{\;AC\;}[]
    5G终端接入时，哪些参数会影响不上报B1事件
    \begin{tasks}(2)
        \task SSB参数
        \task PDCP头压缩配置
        \task 波束类型参数
        \task VRB交织参数
    \end{tasks}
\end{choice}

\begin{choice}{\;ACD\;}[]
    大唐5G基站设备维护时，哪些时钟状态不会造成小区退服
    \begin{tasks}(2)
        \task HoldOver状态
        \task HoldOver超时状态
        \task 预热状态
        \task 锁定状态
    \end{tasks}
\end{choice}

\begin{choice}{\;ABCD\;}[]
    在5G网络中PRACH规划的内容包括哪些
    \begin{tasks}(2)
        \task 频率规划
        \task Preamble格式规划
        \task Preamble配置索引
        \task 根序列索引
    \end{tasks}
\end{choice}

\begin{choice}{\;ABCD\;}[]
    基于数字李生技术和人工智能技术，6G网络将成为具备以下功能的自治网络

    \begin{tasks}(2)
        \task 自配置
        \task 自愈
        \task 自演进
        \task 自生长能力
    \end{tasks}
\end{choice}

\begin{choice}{\;ABD\;}[]
    以下属于GradientBoosting法的选项是
    \begin{tasks}(4)
        \task lightGBM
        \task GBDT
        \task cart树
        \task Xgboost
    \end{tasks}
\end{choice}


\begin{choice}{\;ABCD\;}[]
    移动通信系统中，分集的方式包括
    \begin{tasks}(2)
        \task 极化分集
        \task 天线分集
        \task 时间分集
        \task 频率分集
    \end{tasks}
\end{choice}


\begin{choice}{\;ABCD\;}[]
    NR系统切换过程中，终端上报切换的MR，但是基站未下发切换重配置，可能原因有哪些
    \begin{tasks}(2)
        \task 外部邻区是否配置
        \task 邻小区关系是否配置
        \task 切换判决条件是否合理
        \task 异频是否配置
    \end{tasks}
\end{choice}


\begin{choice}{\;AC\;}[]
    自由空间的传播模型的无线电波损耗与哪些因素有关
    \begin{tasks}(2)
        \task 传播距离
        \task 一般城区、郊区等等地物类型
        \task 电波频率
        \task 天线高度
    \end{tasks}
\end{choice}


\begin{choice}{\;CD\;}[]
    关于5G测量，以下说法正确的是：
    \begin{tasks}(1)
        \task NR测量中只在异频测量中需要gap，同频测量不需要gap
        \task NR测量中测量事件的评估是以beam信号质量进行评估的
        \task NR测量中同一个频率即使子载波间隔不同，也需要配置两个不同的MO
        \task 对于NR测量对象中配置的黑名单列表中的小区，UE不需要进行事件评估和测量
    \end{tasks}
\end{choice}


\begin{choice}{\;AC\;}[]
    NR中上行HARQ采用下列哪种方式
    \begin{tasks}(2)
        \task 异步
        \task 同步
        \task 自适应
        \task 非自适应
    \end{tasks}
\end{choice}


\begin{choice}{\;BC\;}[]
    5G大规模天线中，上行支持的MIMO方式有
    \begin{tasks}(2)
        \task 不支持波束赋型发送
        \task 基于DMRS的SU-MIMO
        \task 支持波束赋型发送
        \task 基于CRS的MU-MIMO
    \end{tasks}
\end{choice}


\begin{choice}{\;ABD\;}[]
    在5GRF优化中关于越区覆盖，以下描述正确的是
    \begin{tasks}(1)
        \task 越区覆盖可能形成“孤岛”，导致掉话
        \task 越区覆盖可能导致UE发射功率受限，引起接入失败
        \task 对于NR系统而言越区覆盖不会对其邻近小区形成干扰
        \task 越区覆盖可能导致越区覆盖基站与周边基站负荷不均。周边基站资源利用率低
    \end{tasks}
\end{choice}


\section{判断题（共30小题，共60分）}

\begin{choice}{\;正确\;}[]
    C-V2X的网络技术包含4G或者5G。
\end{choice}


\begin{choice}{\;错误\;}[]
    5G基本时间单位（Tc）与LTE基本时间单位（Ts）是一样的。
\end{choice}

\begin{choice}{\;正确\;}[]
    5G网络可以根据车联网业务提供高密度连接切片。
\end{choice}

\begin{choice}{\;正确\;}[]
    5G NR系统PUCCH传输的内容包括ACK/NACK。
\end{choice}

\begin{choice}{\;正确\;}[]
    NR系统中，gNB通过RRCrelease消息触发UE进入去激活状态，包含TAC消息。
\end{choice}

\begin{choice}{\;错误\;}[]
    光纤传输采用幅移键控ASK调制方法，即亮度调制，多模光纤性能优于单模光纤。

\end{choice}


\begin{choice}{\;错误\;}[]
    5G的基站功能重构为CU和DU两个功能实体，CU主要处理物理层功能和实时性需求的层 2功能。

\end{choice}

\begin{choice}{\;正确\;}[]
    大规模天线可以是有源天线也可以是无源天线。

\end{choice}

\begin{choice}{\;正确\;}[]

    分离承载是指可以由MeNB或SeNB与GW对接，再通过Xn/Xx接口数据分离。
\end{choice}


\begin{choice}{\;正确\;}[]
    MCGsplitbearer的特性：在双连接里，承载的无线协议在MgNB被分裂，并且承载既属于MCG，也属于SCG。

\end{choice}

\begin{choice}{\;正确\;}[]
    在5G系统中，UE如果携带了TAI是当前注册区的TAI。

\end{choice}


\begin{choice}{\;错误\;}[]
    控制与转发分离：控制面进行统一的策略控制，转发面更关注转发功能的实现。以软件来驱动，实现了软硬件的解耦。

\end{choice}


\begin{choice}{\;错误\;}[]
    移动通信系统中，影响覆盖的菲涅尔区是一个球体。

\end{choice}


\begin{choice}{\;正确\;}[]
    完全成本法按经济用途将成本分为制造成本和非制造成本。

\end{choice}


\begin{choice}{\;正确\;}[]
    移动通信系统中，在上行链路上，扰码是用于基站区分来自不同小区的用户。

\end{choice}


\begin{choice}{\;错误\;}[]
    5G频率的保护带宽必须是对称的。

\end{choice}

\begin{choice}{\;正确\;}[]
    完全成本法的理论依据是：凡是同产品生产有关的消耗都应计入产品成本。

\end{choice}


\begin{choice}{\;正确\;}[]
    边缘计算：适用于实时、短周期数据、本地决策场景。

\end{choice}


\begin{choice}{\;错误\;}[]
    调制指数越高，对设备的性能要求越低，设备费用、复杂性会降低。

\end{choice}


\begin{choice}{\;错误\;}[]
    我国固定电话网目前按照服务对象的网络等级分类，可以分为5级。

\end{choice}


\begin{choice}{\;正确\;}[]
    馈线拐弯应圆滑均匀，弯曲点应尽可能少。

\end{choice}


\begin{choice}{\;正确\;}[]
    目前用在SLAM上的传感器主要分为这两类，一种是基于激光雷达的激光SLAM（Lidar SLAM）和基于视觉的VSLAM（VisualSLAM）。

\end{choice}


\begin{choice}{\;正确\;}[]
    大规模天线阵列正是基于多用户波束成形的原理，在基站端布置多根天线，对几十个目标接收机调制各自的波束。

\end{choice}


\begin{choice}{\;正确\;}[]
    卫星电话手机与普通手机相比，需要额外的卫星通信天线和硬件来支持与卫星通信系统进行通信。
\end{choice}

\begin{choice}{\;正确\;}[]
    在以解决复杂工程技术问题为目标时，追求个别技术指标的先进性可能会导致其他关联技术指标的劣化。

\end{choice}

\begin{choice}{\;错误\;}[]
    为快速对准波束，5G标准采取了分级扫描的策略，即由窄到宽扫描。

\end{choice}

\begin{choice}{\;正确\;}[]
    在密集琥珀区、丘陵城市，及有玻璃幕墙建筑的环境中，选址时要注意信号反射和衍射的影响。

\end{choice}

\begin{choice}{\;正确\;}[]
    人工智能AI在5G+传输网方面可以提升信号处理的准确性和效率。

\end{choice}

\begin{choice}{\;正确\;}[]
    开通双模基站时，需要先进行5G开通调测。

\end{choice}

\begin{choice}{\;错误\;}[]
    5G所有场景都支持自包含帧结构。

\end{choice}








